% from ega4-2-fr.tex
\oldpage[IV-2]{88}

des $x\in X$ tels que $f(x)=0$ est de codimension $\leq1$ dans $X$; elle est
de codimension $1$ si $Z$ ne contient aucune composante irréductible de $X$.
\end{corollary}

\begin{proof}
On peut se borner au cas où $\mathcal L=\mathcal O_X$. Si $y$ est un point
générique d'une composante irréductible $Y$ de $Z$, l'idéal $(f_y)$ de
$\mathcal O_{X,y}$ doit être tel que $\mathcal O_{X,y}/(f_y)$ n'ait qu'un
seul idéal premier, ce qui signifie que $f_y$ engendre un idéal de définition
de l'anneau local noethérien $\mathcal O_{X,y}$; on a donc
$\operatorname{codim}(Y,X)\leq1$ (5.1.7); si $Z$ ne contient aucune
composante irréductible de $X$, on ne peut avoir
$\operatorname{codim}(Y,X)=0$ en vertu de $(0,14.2.1)$.
\end{proof}

\begin{proposition}[5.1.9]
\label{IV.2.5.1.9-fr}
Soient $X$ un préschéma localement noethérien, $Y$ une partie fermée de $X$.
Si $x\in Y$ est tel que $\mathcal O_{X,x}$ soit un anneau caténaire, on a
\begin{align}
\operatorname{codim}_x(Y,X)
&=\dim(\mathcal O_{X,x})-\operatorname{codim}(\overline{\{x\}},Y)\notag\\
&=\dim(\mathcal O_{X,x})-\dim(\mathcal O_{Y,x}).
\tag{5.1.9.1}\label{IV.2.5.1.9.1-fr}
\end{align}
\end{proposition}

\begin{proof}
Soient $Y_i$ ($1\leq i\leq n$) les composantes irréductibles de $Y$
contenant $x$ (qui sont en nombre fini puisque $Y$ est localement
noethérien), et soit $y_i$ le point générique de $Y_i$. Si l'on pose
$A=\mathcal O_{X,x}$, et si $\mathfrak p_i$ est l'idéal premier de $A$
correspondant à $Y_i\cap\operatorname{Spec}(A)$, les parties fermées
irréductibles de $Y$ contenant $x$ correspondent biunivoquement aux idéaux
premiers de $A$ contenant un des $\mathfrak p_i$, donc
$\dim(\mathcal O_{Y,x})=\sup_i\dim(A/\mathfrak p_i)$; d'autre part, on a
$\mathcal O_{X,y_i}=A_{\mathfrak p_i}$, et l'hypothèse sur $A$ entraîne
$\dim(A)=\dim(A/\mathfrak p_i)+\dim(A_{\mathfrak p_i})$ $(0,16.1.4)$; la
conclusion résulte donc de la relation
$\operatorname{codim}_x(Y,X)=\inf_i(\dim(\mathcal O_{X,y_i}))$ (5.1.2 et
$(0,14.2.6)$).
\end{proof}

\begin{proposition}[5.1.10]
\label{IV.2.5.1.10-fr}
\begin{enumerate}
\item[(i)] Dans tout préschéma $X$, toute partie localement constructible non
vide $E$ contient un point $x$ tel que $\{x\}$ soit une partie localement
fermée de $X$ (ou, ce qui revient au même, tel que $x$ soit isolé dans
$\overline{\{x\}}$).
\item[(ii)] Soient $X$ un préschéma localement noethérien, $x$ un point de
$X$ tel que $\{x\}$ soit localement fermé dans $X$; alors on a
$\dim(\overline{\{x\}})\leq1$; tout point $y\ne x$ de
$\overline{\{x\}}$ est par suite fermé dans $X$.
\end{enumerate}
\end{proposition}

\begin{proof}
(i) Le résultat est un cas particulier du lemme suivant:
\end{proof}

\begin{lemma}[5.1.10.1]
\label{IV.2.5.1.10.1-fr}
Soit $X$ un espace topologique ayant la propriété suivante: pour toute partie
localement fermée $Z\ne\varnothing$ de $X$, il existe une partie $Z'$ de $Z$,
localement fermée dans $X$ (ou dans $Z$, ce qui revient au même) et un point
$x\in Z'$, fermé dans $Z'$. Alors toute partie localement fermée
$Z\ne\varnothing$ de $X$ contient un point $x$ isolé dans
$\overline{\{x\}}$.
\end{lemma}

\begin{proof}
En effet, soit $Z'\subset Z$ une partie localement fermée de $Z$ contenant un
point $x$ tel que $x$ soit fermé \emph{dans} $Z'$. Il y a un voisinage ouvert
$U$ de $x$ dans $X$ tel que $Z'\cap U$ soit fermé dans $U$, donc $x$ est aussi
fermé dans $U$; cela signifie que $U\cap\overline{\{x\}}=\{x\}$, autrement
dit que $x$ est isolé dans $\overline{\{x\}}$.

Le lemme s'applique à tout espace sous-jacent à un préschéma $X$, car $Z$ est
alors aussi espace sous-jacent à un préschéma (I, 5.2.1) et il suffit de
prendre pour $Z'$ un ouvert affine dans $Z$, qui est un espace de Kolmogoroff
quasi-compact, donc contient un point fermé $(0_{\mathrm I},2.1.3)$.
\end{proof}

